% =============================================================================
% 5. fejezet — A rendszer tervezése és megvalósítása
% =============================================================================

Ez a fejezet a DermAI rendszer mérnöki megvalósítását mutatja be: a betanított
osztályozó modellt egy teljes, végponttól végpontig működő alkalmazássá
egészíti ki, amely okostelefonon és webböngészőben egyaránt elérhető. Először a
követelményeket és a rendszerarchitektúrát ismertetjük, majd komponensenként a
megvalósítást: a modellt kiszolgáló réteget, a backendet, a mobil- és
webalkalmazást, a hitelesítést és a felhőalapú adattárolást, végül a hálózati
hibatűrést és a biztonsági megfontolásokat.

% =============================================================================
\section{Követelmények}
% =============================================================================

A rendszerrel szemben támasztott elvárások két forrásból fakadnak: egyrészt a
dolgozat céljából -- egy nem szakértő felhasználó által, hétköznapi eszközön
végigvihető szűrési folyamatból --, másrészt a megvalósítás valós
körülményeiből: a változó hálózati elérhetőségből, az érzékeny egészségügyi
adatok kezeléséből és a korlátozott, gyakran cserélődő GPU-erőforrásból. Az
alábbiakban előbb a funkcionális, majd a nem-funkcionális elvárásokat foglaljuk
össze.

\subsection{Funkcionális követelmények}

A rendszerrel szemben támasztott fő funkcionális elvárások a következők:

\begin{itemize}
  \item A felhasználó fényképet készíthet vagy tölthet fel egy
  bőrelváltozásról (kamera vagy galéria).
  \item A rendszer a feldolgozás előtt ellenőrzi a kép minőségét és
  relevanciáját, és elutasítja az értékelhetetlen felvételeket.
  \item A bőrelváltozást egy ViT-alapú modell osztályozza
  (jó\-indulatú / rosszindulatú), konfidencia-értékkel.
  \item A rendszer természetes nyelvű magyarázatot ad az eredményről, az
  ABCDE-szabály figyelembevételével.
  \item A felhasználó utólagos kérdéseket tehet fel az elváltozással
  kapcsolatban (párbeszéd).
  \item A felhasználó bejelentkezik, és a feltöltött képek, valamint a
  beszélgetések a saját fiókjához kötve, biztonságosan tárolódnak.
\end{itemize}

A rendszer fő használati eseteit a \ref{fig:usecase}.~ábra foglalja össze.

\begin{figure}[H]
  \centering
  \includegraphics[width=\textwidth]{uml/usecase.png}
  \caption{A rendszer használati eset diagramja: a felhasználó fő céljai, valamint
    az ,,elváltozás elemzése'' use case kötelező lépései (képvalidáció,
    osztályozás, magyarázat).}
  \label{fig:usecase}
\end{figure}

\subsection{Nem-funkcionális követelmények}

\begin{itemize}
  \item \textbf{Reszponzivitás:} az AI-válaszok token-szintű, fokozatos
  megjelenítése, hogy a felhasználó ne egy hosszú várakozás után, hanem
  folyamatosan kapjon visszajelzést.
  \item \textbf{Hibatűrés:} a rendszer alkalmazkodjon a változó hálózati
  körülményekhez és a modellt kiszolgáló komponens időszakos
  elérhetetlenségéhez.
  \item \textbf{Biztonság:} a felhasználói adatok hitelesített hozzáférés mögött,
  felhasználónként elkülönítve tárolódjanak; a titkos kulcsok ne kerüljenek a
  forráskódba.
  \item \textbf{Hordozhatóság:} a modellt kiszolgáló komponens címe futás közben
  módosítható legyen, hogy a (gyakran változó címen elérhető) GPU-háttér
  újraindítása ne igényelje az alkalmazás újratelepítését.
\end{itemize}

Ezek a nem-funkcionális elvárások közvetlenül alakították a rendszer felépítését:
a reszponzivitást a token-szintű adatfolyam (SSE), a hibatűrést az
újrapróbálkozó hálózati réteg és az önjavító modellkliens, a hordozhatóságot a
modellkiszolgáló futásidejű címfrissítése, a biztonságot pedig a hitelesített
végpontok és a felhasználónként elkülönített tárolás valósítja meg. E megoldások
részleteit a következő szakaszok tárgyalják.

% =============================================================================
\section{Rendszerarchitektúra}
% =============================================================================

A DermAI több, jól elkülönülő komponensből álló, elosztott rendszer. A két
kliens (Android alkalmazás és webes felület) külön-külön backendhez kapcsolódik,
amelyek azonban közös elveken működnek, és ugyanazt a modellt kiszolgáló réteget,
illetve felhőszolgáltatásokat használják. A komponensek és kapcsolataik
áttekintését a \ref{fig:architecture}.~ábra mutatja.

\begin{itemize}
  \item \textbf{Android alkalmazás} (Kotlin) -- a mobil kliens; a
  \textbf{Flask} alapú backendhez \cite{flask} kapcsolódik.
  \item \textbf{Webes felület} (HTML/JavaScript) -- a böngészőkliens; a
  \textbf{FastAPI} alapú backendhez \cite{fastapi} kapcsolódik.
  \item \textbf{Modellt kiszolgáló réteg} -- egy MCP-szerver
  \cite{mcp2024}, amely a GPU-n futó ViT modellt szólítja meg.
  \item \textbf{Gemini 2.5 Flash} \cite{gemini2023} -- a képvalidációt, a
  magyarázatgenerálást és a párbeszédet végző nagy nyelvi modell.
  \item \textbf{Firebase} \cite{firebase} -- hitelesítés (Authentication),
  adatbázis (Firestore) és fájltárolás (Storage).
\end{itemize}

\begin{figure}[H]
  \centering
  \includegraphics[width=\textwidth]{uml/component.png}
  \caption{A rendszer komponens-diagramja: a kliensek (Android, web), a backendek
    (Flask, FastAPI), a modellt kiszolgáló MCP-réteg és a GPU-n futó ViT, valamint
    a Firebase és a Gemini szolgáltatások, a köztük lévő kommunikációs
    protokollokkal.}
  \label{fig:architecture}
\end{figure}

A rendszer rétegzett, a felelősségi körök szerint szétválasztott felépítést
követ. A kliensek kizárólag a felhasználói interakcióért és a megjelenítésért
felelnek; az üzleti logikát -- a kérések összehangolását, a külső szolgáltatások
hívását és a hozzáférés-ellenőrzést -- a backend végzi; a számításigényes
modellinferencia pedig egy külön, GPU-val rendelkező rétegbe került. Ez a
szétválasztás lehetővé teszi, hogy az egyes komponensek egymástól függetlenül
legyenek fejleszthetők, cserélhetők és skálázhatók.

A backend tölti be a központi összehangoló (orkesztrátor) szerepét: ez fogadja a
kliensek kéréseit, ellenőrzi a hitelesítő tokent, majd sorrendben meghívja a
képvalidációt (Gemini), az osztályozást (a modellt kiszolgáló rétegen keresztül a
ViT) és a magyarázatgenerálást, közben pedig gondoskodik az adatok tartós
tárolásáról (Firebase). A kliensek így soha nem érik el közvetlenül sem a
modellt, sem a nagy nyelvi modellt -- minden külső hívás a backenden,
hitelesített végpontok mögött történik, ami biztonsági szempontból is előnyös.

Két tervezési döntés emelendő ki. Egyrészt a modellt kiszolgáló réteg tudatosan
le van választva a backendről: a ViT egy önálló, GPU-s környezetben fut, amelyet
a backend az MCP-rétegen keresztül ér el; a GPU-háttér (gyakran változó) nyilvános
címe futás közben átállítható, így annak újraindítása vagy cseréje nem igényli a
backend módosítását.
Másrészt a két mesterséges intelligencia komponens szerepe szigorúan elkülönül: a
kvantitatív döntést a célzottan betanított ViT hozza, míg a nagy nyelvi modell
(Gemini) szerepe a bemenet ellenőrzésére, az eredmény közérthető magyarázatára és
a párbeszédre korlátozódik.

\subsection{Egy elemzési kérés életciklusa}

Egy kép elemzésének tipikus folyamata a következő lépésekből áll:

\begin{enumerate}
  \item A kliens elküldi a képet és a hitelesítő tokent a backendnek.
  \item A backend ellenőrzi a tokent, majd a képet a Gemini modellnek továbbítja
  \textit{minőség- és relevanciaellenőrzésre}. Ha a kép nem értékelhető (pl.
  nem bőrelváltozás, túl távoli, életlen vagy több léziót tartalmaz), a folyamat
  itt megszakad, és a felhasználó célzott visszajelzést kap.
  \item Megfelelő kép esetén a backend feltölti a képet a felhasználó
  tárhelyére (Firebase Storage), majd az MCP-rétegen keresztül lefuttatja a
  ViT-osztályozást.
  \item A backend létrehozza vagy frissíti a munkamenet rekordját a Firestore
  adatbázisban.
  \item A backend a klasszifikációs eredményt belefoglalja a Gemini számára
  összeállított utasításba, és a magyarázatot \textbf{token-szintű
  adatfolyamként} (SSE) küldi vissza a kliensnek.
\end{enumerate}

A folyamat teljes lefutását a \ref{fig:sequence-analyze}.~ábra szekvenciadiagramja
szemlélteti.

\begin{figure}[H]
  \centering
  \includegraphics[width=\textwidth]{uml/sequence_analyze.png}
  \caption{Egy elemzési kérés (\texttt{/analyze}) szekvenciadiagramja: a kliens, a
    backend, a Gemini, az MCP-szerver, a ViT és a Firebase interakciói. Az elágazás
    (\textit{alt}) az értékelhetetlen kép esetét, a ciklus (\textit{loop}) a
    token-szintű (SSE) válasz továbbítását jelöli.}
  \label{fig:sequence-analyze}
\end{figure}

A folyamat két eleme érdemel külön figyelmet. A Gemini-alapú \emph{előzetes
képvalidáció} kapuként működik: az értékelhetetlen felvételeket még a
számításigényesebb osztályozás előtt kiszűri, csökkentve a felesleges
erőforrás-használatot és annak kockázatát, hogy a felhasználó értelmezhetetlen
képre kapjon megtévesztő eredményt. A magyarázat \emph{token-szintű,
szerver-küldött eseményekkel} (SSE) történő visszaküldése pedig azt biztosítja,
hogy a felhasználó ne egy hosszú várakozás után, hanem fokozatosan, a
generálással egyidejűleg kapja meg a választ.

% =============================================================================
\section{A modellt kiszolgáló komponens: MCP-szerver}
% =============================================================================

A betanított ViT modell a backendtől elkülönített, GPU-val rendelkező
környezetben (Google Colab) fut, és HTTP-végponton fogadja a
klasszifikációs kéréseket. A backend és e modellkiszolgáló közötti
kommunikációt egy MCP- (Model Context Protocol) szerver \cite{mcp2024} közvetíti,
amely a backendből indított alfolyamatként, szabványos be- és kimeneti (stdio)
csatornán, JSON-RPC 2.0 üzenetekkel \cite{jsonrpc2013} kommunikál. A szerver
egyetlen eszközt (tool) tesz közzé, a \texttt{classify\_lesion} függvényt; egy
klasszifikációs kérés üzenetformátuma például:

\begin{verbatim}
{"jsonrpc": "2.0", "id": 1, "method": "tools/call",
 "params": {"name": "classify_lesion",
            "arguments": {"image_base64": "..."}}}
\end{verbatim}

Az MCP-szerver a kapott, base64 kódolású képet továbbítja a GPU-n futó modell
\texttt{/classify} végpontjának, és a választ -- a címkét, a konfidenciát és az
osztályonkénti pontszámokat -- adja vissza a backendnek. A megoldás egyik
kulcseleme egy \texttt{update\_model\_url} metódus, amely lehetővé teszi a
modellkiszolgáló címének \textbf{futás közbeni} módosítását: mivel a GPU-háttér
újraindításakor a nyilvános cím (alagút) megváltozik, e mechanizmus révén a
backend új cím esetén sem igényel újraindítást. A modellt kiszolgáló réteg
ilyen módon történő leválasztása lehetővé teszi, hogy a GPU-háttér bármikor
újraindítható vagy lecserélhető legyen anélkül, hogy a backend kódját módosítani
kellene.

% =============================================================================
\section{A backend}
% =============================================================================

A rendszer két backendet tartalmaz, amelyek a két klienst szolgálják ki, de
azonos belső felépítést követnek. A mobil kliens egy \textbf{Flask}
\cite{flask} alapú szervert, a webes felület pedig egy \textbf{FastAPI}
\cite{fastapi} alapú szervert használ. Mindkét backend ugyanazt az MCP-klienst,
ugyanazt a Gemini-integrációt és ugyanazokat a Firebase-szolgáltatásokat
használja. A FastAPI aszinkron működése a webes streaminghez, a Flask
egyszerűsége pedig a mobil végpontok kiszolgálásához bizonyult kényelmesnek.

A két szerver egy közös belső magra épül: ugyanazokat a modulokat használják a
modellkiszolgáló elérésére (MCP-kliens), a Gemini-integrációra és a
Firebase-műveletekre. Így a kétféle kiszolgálás viselkedése következetes marad,
és az üzleti logika nem duplikálódik a két kódbázisban -- a két backend
gyakorlatilag csak a webes keretrendszerben és a kérés-/válaszkezelés módjában
(szinkron, illetve aszinkron) tér el.

\subsection{Végpontok}

A webes (FastAPI) backend főbb végpontjait a \ref{tab:endpoints}.~táblázat
foglalja össze; a mobil (Flask) backend ezekkel funkcionálisan egyenértékű
végpontokat biztosít.

\begin{table}[H]
  \centering
  \caption{A webes backend (FastAPI) főbb végpontjai.}
  \label{tab:endpoints}
  \begin{tabular}{lll}
    \hline
    \textbf{Metódus} & \textbf{Útvonal} & \textbf{Funkció} \\
    \hline
    POST & \texttt{/analyze}            & Kép elemzése (validálás, osztályozás, magyarázat) \\
    POST & \texttt{/chat}               & Utólagos kérdés (kép nélkül), párbeszéd \\
    GET  & \texttt{/history/\{id\}}     & Egy korábbi munkamenet lekérése \\
    GET  & \texttt{/sessions}           & A felhasználó munkameneteinek listája \\
    PUT  & \texttt{/model-url}          & A modellkiszolgáló címének frissítése \\
    GET  & \texttt{/health}             & Állapotellenőrzés \\
    \hline
  \end{tabular}
\end{table}

A végpontok REST-elvűek és állapotmentesek: minden kérés önállóan hordozza a
hitelesítő tokent, a munkamenet állapotát pedig nem a szerver memóriája, hanem a
Firestore tárolja. Ez egyszerűsíti a kiszolgálást, és megkönnyíti a backend
esetleges vízszintes skálázását.

\subsection{A Gemini integrációja és a képvalidáció}

A nagy nyelvi modell (Gemini 2.5 Flash) a rendszerben három feladatot lát el. Az
\textbf{elsődleges szűrés} során a feltöltött képet egy rögzített utasítással
értékeli, amely egyetlen kódot kér vissza: \texttt{OK} (értékelhető),
\texttt{NOT\_SKIN}, \texttt{TOO\_FAR}, \texttt{MULTIPLE}, \texttt{BLURRY} vagy
\texttt{NOT\_ASSESSABLE}. Ha a kód nem \texttt{OK}, a felhasználó a hibának
megfelelő, célzott útmutatást kap (például: tartsa közelebb a kamerát, vagy
fényképezzen egyetlen léziót). A \textbf{magyarázatgenerálás} során a Gemini a
ViT klasszifikációs eredményét beépítve, az ABCDE-szabály szerint értelmezi az
elváltozást, közérthető nyelven. Végül a Gemini kezeli a felhasználó
\textbf{utólagos kérdéseit} is, a beszélgetés előzményeinek figyelembevételével.
Ez a munkamegosztás tudatos tervezési döntés: a kvantitatív osztályozást a
célzottan betanított ViT végzi, a nyelvi modell pedig az erősségeire -- a
közérthető magyarázatra és a párbeszédre -- összpontosít.

\subsection{Token-szintű adatfolyam (SSE)}

A reszponzív felhasználói élmény érdekében a Gemini válaszait nem egyben, hanem
token-szintű adatfolyamként továbbítjuk a kliensnek, szerver-küldött események
(Server-Sent Events) formájában \cite{whatwg_sse}. Az adatfolyam egyes
eseményei egy-egy szövegdarabot tartalmaznak, a lezárást pedig egy különleges
jelző jelzi:

\begin{verbatim}
data: {"session_id": "...", "image_url": "...", "classification": {...}}
data: {"token": "Az elváltozás "}
data: {"token": "aszimmetrikus, ..."}
data: [DONE]
\end{verbatim}

Így a felhasználó a szöveget a generálással egyidejűleg, folyamatosan látja
megjelenni. Az első adatfolyam-esemény a munkamenet metaadatait (azonosító,
feltöltött kép URL-je, klasszifikáció) hordozza, így a kliens már a szöveges
válasz előtt frissítheti a felületet.

% =============================================================================
\section{A mobilalkalmazás}
% =============================================================================

A mobil kliens natív Android alkalmazás, Kotlin nyelven, az MVVM (Model--View--ViewModel)
mintát követve. Indításkor egy bejelentkező képernyő (\texttt{LoginActivity})
fogadja a felhasználót, ahol Google-fiókkal jelentkezhet be; sikeres hitelesítés
után a fő képernyőre (\texttt{MainActivity}) jut.

A fő képernyő chat-alapú felületet biztosít: a felhasználó a kamerával
fényképet készíthet, vagy a galériából választhat képet, majd szöveges üzenettel
együtt elküldheti azt elemzésre. A beérkező AI-válaszokat a felület token-szintű
adatfolyamként, folyamatosan jeleníti meg, a klasszifikáció eredményét pedig egy
színkódolt jelvény (jó\-indulatú / rosszindulatú) emeli ki. A jobb felső sarokban
elérhető menü három funkciót kínál: új beszélgetés indítása (\textit{New Chat}),
az alkalmazás ismertetője (\textit{About}) és a kijelentkezés (\textit{Sign out}).

Az MVVM-minta szétválasztja a felületet és az állapotkezelést, ami egy
adatfolyam-alapú, hosszan futó interakciónál -- a token-szintű válaszok
fokozatos megjelenítésénél -- különösen előnyös. A felület a klasszifikáció
eredményét tudatosan \emph{kockázati jelzésként}, nem diagnózisként jeleníti
meg, összhangban a rendszer tájékoztató jellegével.

\subsection{Kameramodul és makrófelvétel}

A bőrelváltozások fényképezéséhez a közeli (makró) fókusz elengedhetetlen. Az
alkalmazás ezért egy dedikált \texttt{CameraActivity} osztályt használ, amely a
CameraX könyvtárra épül. A modul a \texttt{CameraManager} segítségével
felsorolja az eszköz hátsó kameráit, és a legkisebb fókusztávolságot támogató
kamerát választja ki, lehetővé téve az elváltozás éles, közeli felvételét; a kép
rögzítése a legjobb minőséget célzó módban történik. Így az alkalmazás a makró
lencsével rendelkező készülékeken is kihasználja azok közelfókuszáló képességét.

\subsection{Hálózati réteg}

Az alkalmazás hálózati rétege az OkHttp és a Retrofit könyvtárakra épül. A
streaming (SSE) válaszokat egy Kotlin \texttt{Flow} kezeli, az állapotot pedig
egy \texttt{ChatViewModel} tartja karban, így az adatfolyam a képernyő-forgatást
is túléli. Minden kéréshez egy elfogó (\texttt{AuthInterceptor}) automatikusan
hozzáfűzi a felhasználó Firebase azonosítótokenjét. A backend címét egy
beállítási réteg (\texttt{NetworkConfig}) tárolja a \texttt{SharedPreferences}-ben,
amely a fordításkor beállított alapértelmezett címről indul -- ez teszi lehetővé,
hogy az alkalmazás ott keresse a backendet, ahol az ténylegesen elérhető.

% =============================================================================
\section{A webes felület}
% =============================================================================

A webes felület egy egyoldalas, sötét témájú alkalmazás, amely a FastAPI backend
által kiszolgált sablonra épül. A felület a Firebase JavaScript SDK-t használja a
Google-alapú bejelentkezéshez: amíg a felhasználó nincs bejelentkezve, egy
bejelentkező réteg takarja a felületet, sikeres hitelesítés után pedig
megjelenik a chat-felület. Minden, a backendhez intézett kérés tartalmazza a
felhasználó hitelesítő tokenjét. A felületen egy beépített, lapváltás nélküli
ismertető (About) ablak érhető el, amely megőrzi a beszélgetés addigi tartalmát;
kijelentkezéskor a felület a helyi állapotot (üzenetek, munkamenet-azonosító) is
törli. A token-szintű válaszokat a böngésző a szerver-küldött események
feldolgozásával jeleníti meg fokozatosan.

A webes kliens szándékosan keretrendszer nélküli (sima HTML, CSS és JavaScript):
ez az egyszerű, build-lépést nem igénylő felépítés elegendő a chat-alapú
felülethez, és a FastAPI által kiszolgált egyetlen sablonból közvetlenül
futtatható. Funkcionálisan a webes és a mobil kliens ugyanazt a folyamatot
kínálja, így a felhasználó a platformtól függetlenül azonos élményt kap.

% =============================================================================
\section{Hitelesítés és felhőalapú adattárolás}
% =============================================================================

A rendszer a Firebase szolgáltatásait \cite{firebase} használja a
hitelesítéshez és az adattároláshoz.

\subsection{Hitelesítés}

A bejelentkezés mindkét platformon Google-fiókkal, a Firebase Authentication
segítségével történik. A kliens a sikeres bejelentkezés után egy
azonosító\-tokent (ID token) kap, amelyet minden backend-kéréshez mellékel az
\texttt{Authorization} fejlécben. A backend e tokent szerveroldalon ellenőrzi a
Firebase Admin SDK segítségével, és a kérésből kinyeri a felhasználó egyedi
azonosítóját (\texttt{uid}) és e-mail-címét; érvénytelen vagy hiányzó token
esetén a kérést elutasítja.

\subsection{Adatbázis és fájltárolás}

A munkameneteket a Firestore adatbázis tárolja, munkamenetenként egy
dokumentumban, amely tartalmazza a tulajdonos azonosítóját és e-mail-címét, a
klasszifikáció eredményét, a feltöltött kép azonosítóját és a beszélgetés
üzeneteit. A dokumentum szerkezete vázlatosan:

\begin{verbatim}
sessions/{session_id}
   user_id        : a tulajdonos Firebase-azonosítója
   user_email     : a tulajdonos e-mail-címe
   created_at     : létrehozás időbélyege
   source         : "web" vagy "android"
   image_id       : a feltöltött kép azonosítója
   classification : {label, confidence, all_scores}
   messages       : [{role, text, timestamp}, ...]
\end{verbatim}

A feltöltött képeket a Firebase Storage tárolja, \textbf{felhasználónként
elkülönített} útvonalon: \texttt{users/\{uid\}/lesions/\{image\_id\}.jpg}. Így a
képek a feltöltő felhasználó fiókjához kötődnek, és a munkamenetek listázásakor a
backend a tulajdonos azonosítója szerint szűr, biztosítva, hogy minden
felhasználó kizárólag a saját adatait érje el.

% =============================================================================
\section{Hálózati robusztusság és hibatűrés}
% =============================================================================

Mivel a rendszer több, egymástól független szolgáltatásra támaszkodik -- és a
modellt kiszolgáló GPU-háttér címe is változhat --, a hálózati hibatűrés
kiemelt tervezési szempont volt. A megvalósított mechanizmusok:

\begin{itemize}
  \item \textbf{Újrapróbálkozás exponenciális visszalépéssel.} A REST-hívásokhoz
  egy OkHttp-elfogó (interceptor), az SSE-adatfolyamhoz pedig a Kotlin
  \texttt{Flow} \texttt{retryWhen} operátora biztosít automatikus
  újrapróbálkozást átmeneti hálózati hiba esetén.
  \item \textbf{Az MCP-kliens önjavítása.} A backend MCP-kliense a sikertelen
  hívásokat többször újrapróbálja, és szükség esetén újraindítja a modellt
  kiszolgáló alfolyamatot.
  \item \textbf{Dinamikus címkezelés.} A modellkiszolgáló címe futás közben
  frissíthető (\texttt{update\_\allowbreak model\_\allowbreak url}, illetve a \texttt{/model-url}
  végpont), a backend címe pedig a kliensoldali beállításokban módosítható, így
  a rendszer alkalmazkodik a változó nyilvános címekhez.
\end{itemize}

% =============================================================================
\section{Biztonsági megfontolások}
% =============================================================================

A rendszer a titkos kulcsokat (a Gemini API-kulcsot és a Firebase hozzáférési
adatokat) környezeti változókból, futásidőben tölti be, és \textbf{soha nem
tárolja a forráskódban}; e fájlokat a verziókezelés is kizárja. A backend minden
védett végponton ellenőrzi a Firebase azonosítótokent, így a klasszifikációs és
párbeszéd-funkciók kizárólag hitelesített felhasználók számára érhetők el. A
felhasználói adatok -- a feltöltött képek és a munkamenetek -- felhasználónként
elkülönítve tárolódnak, így egy felhasználó nem férhet hozzá más adataihoz.

Hangsúlyozandó továbbá, hogy a rendszer kialakítása tudatosan
\textbf{tájékoztató jellegű}: a klasszifikációt egy célzottan betanított
képi osztályozó modell végzi, míg a nagy nyelvi modell szerepe a magyarázatra és a
párbeszédre korlátozódik. A rendszer nem helyettesíti a szakorvosi diagnózist,
hanem szűrőeszközként segíti a felhasználót a dermatológiai konzultáció
szükségességének megítélésében.
